\scanpage{73}
considéré comme sous-groupe du groupe de Brauer formé des éléments annulés par 2. Donc \uline{la \unsure{structure|strcuture|struuture} de} $Q(A)$ \uline{est connue quand on connait les groupes} $H^{\handwritten{i}}(A,\underline{Z}/2\underline{Z})$ \uline{pour} $i=1,2$, \uline{soient} $D$ \uline{et} $B$, \uline{et le cup-produit} $DxD\handwritten{\to}B$. L'énoncé précédent montre \deleted{en effet}\inserted{\handwritten{même}} que le noyau de l'homomorphisme surjectif (\handwritten{*}) est aussi le sous-groupe engendré par les éléments
\[
e_{a+u+v}-e_{a+u}-e_{a+v}+e_a
\]
pour les éléments $a,u,v$, dans $D$ tels que $uv=0$.

On en conclut en particulier que le gradué associé à $Q(A)$ est isomorphe au quotient de l'algèbre symétrique de $D$ sur $\underline{Z}$, par \deleted{\unsure{les|un}} l'idéal engendré par les éléments d'ordre 2 de la forme $ab-a'b'$, où $a,b,a',b'$, sont comme ci-dessus, ou ce qui revient au, même, engendré par les éléments $uv$ de $\mathrm{Symm}^{\handwritten{2}}(D)$ pour $u,v\quad D$ tels que $u.v=0$ dans $B^{\deleted{\unsure{x}}}$. (N.B. il faudrait éclaircir si $\mathrm{grad}^{\handwritten{2}}Q(A)$ \deleted{est} $\handwritten{\to}B$ est \deleted{\unsure{bien|très}}\inserted{\handwritten{\unsure{tjs|tjrs}}} injectif, en d'autres termes si le noyau de $\mathrm{Symm}^{\handwritten{2}}(D)\quad B$ est engendré par des éléments de la forme \deleted{\unsure{uv}} $uv$).

On conclut du \deleted{\unsure{ceci|tout}} théorème de structure que si ${\unsure{\mathbb{H}|H}}^{\handwritten{\unsure{\circ\cdot|\bullet}}}$ est un anneau gradué à degrés positifs, alors la donnée d'un homomorphisme \deleted{additif}\inserted{\handwritten{d'anneaux augmentés}} de $Q(A)$ dans $\handwritten{\mathbb{Z}\times(1+}\overset{\handwritten{\scriptscriptstyle\wedge}}{\mathbb{H}}^{\handwritten{+}}\handwritten{)}$ équivaut à la donnée d'un homomorphisme
\[
w^{\handwritten{1}}:D\handwritten{\longrightarrow}\mathbb{H}^{\handwritten{1}}
\]
tel que $W^{\handwritten{1}}(a)W^{\handwritten{1}}(b)=W^{\handwritten{1}}(a')W^{\handwritten{1}}(b')$ lorsque $a,b,a',b'$ sont comme ci-dessus, ou ce qui revient au même, tel que $W^{\handwritten{1}}(u)W^{\handwritten{1}}(v)=0$ si $u.v=0$ dans $B$; \deleted{\unsure{on|en}} d'ailleurs $\mathrm{grad}\ Q(A)$ joue un rôle universel pour ses propriétés, i.e. la donnée de $W$ équivaut aussi à la donnée d'un homomorphisme \deleted{addi} d'anneaux gradués
\[
\mathrm{grad}\ Q(A)\handwritten{\to}H^{\handwritten{\bullet}}\ .
\]
Le choix le plus naturel pour $H^{\handwritten{\bullet}}$, du point de vue ``topologique'', semble
\[
H_{\deleted{\unsure{i|1}}}^{\handwritten{\bullet}}=H^{\handwritten{\bullet}}(A,\underline{Z}/2\underline{Z})\ ,
\]
muni du cup-produit, ce qui nous donne des classes de \uline{Stiefel-Whitney}
\[
W^{\handwritten{i}}:Q(A)\handwritten{\to}H^{\handwritten{i}}(A,\underline{Z}/2\underline{Z})\ .
\]
\unsure{Serre|Berre} suggère de \deleted{prendre} considérer aussi une définition plus fine des classes de Stiefel-Whitney, en remplaçant \deleted{\unsure{le grad|le grou}} l'extension séparable algébrique maximale de $A$ par la sous-extension réunion des sous-extensions finies de rang une puissance de 2; cela donnerait déjà un $H^{\handwritten{2}}$ ``plus fin'', i.e. s'envoyant dans le $H^{\handwritten{2}}$ habituel ...
% archive sheet number: 72
